%%%%%%%%%%%%%%%%%%%%%%%%%%%%%%%%%%%%%%%%%%%%%%%%%%%%%%%%%%%%%%%%%%%%%%
% Lớp 8 | Chương I. Đa thức | Đơn thức
% Dạng 8D1?1-3: Tính giá trị của đơn thức
%%%%%%%%%%%%%%%%%%%%%%%%%%%%%%%%%%%%%%%%%%%%%%%%%%%%%%%%%%%%%%%%%%%%%%


% nguồn: De_AI_Toan8_cuoi_hoc_ki_i_1 (nhập 21/11/2026)
\begin{bt}%[8D1N1-3]%[ma:c7eef580]
	Tính giá trị của mỗi đơn thức sau:
a) $A = -3x^2y^3$ tại $x = 2$, $y = -1$;
b) $B = \dfrac{1}{2}x^3y^2$ tại $x = -2$, $y = 3$.
	\loigiai{
		a) Thay $x = 2$, $y = -1$ vào $A$, ta được $A = -3 \cdot 2^2 \cdot (-1)^3 = -3 \cdot 4 \cdot (-1) = 12$. \\
b) Thay $x = -2$, $y = 3$ vào $B$, ta được $B = \dfrac{1}{2} \cdot (-2)^3 \cdot 3^2 = \dfrac{1}{2} \cdot (-8) \cdot 9 = -36$.
	}
\end{bt}

% nguồn: De_AI_Toan8_cuoi_hoc_ki_i_1 (nhập 21/11/2026)
\begin{bt}%[8D1N1-3]%[ma:1295ea5b]
	Cho đơn thức $P = 4x^2y \cdot \left(-\dfrac{1}{2}xy^2\right)$.
a) Thu gọn đơn thức $P$.
b) Tính giá trị của $P$ tại $x = -1$, $y = 2$.
	\loigiai{
		a) $P = 4 \cdot \left(-\dfrac{1}{2}\right) \cdot (x^2 \cdot x) \cdot (y \cdot y^2) = -2x^3y^3$. \\
b) Thay $x = -1$, $y = 2$ vào $P$, ta được $P = -2 \cdot (-1)^3 \cdot 2^3 = -2 \cdot (-1) \cdot 8 = 16$.
	}
\end{bt}

% nguồn: De_AI_Toan8_cuoi_hoc_ki_i_1 (nhập 21/11/2026)
% nguồn: De_AI_Toan8_cuoi_hoc_ki_i_1 (nhập 21/11/2026)
\begin{bt}%[8D1N1-3]%[ma:cdf16e83]
	Một hình lập phương có độ dài cạnh là $2x$ (cm) với $x > 0$.
a) Viết đơn thức biểu thị thể tích của hình lập phương đó và cho biết bậc của đơn thức.
b) Tính thể tích của hình lập phương khi $x = 1{,}5$.
	\loigiai{
		a) Thể tích hình lập phương là $V = (2x)^3 = 8x^3$ ($\text{cm}^3$). Đơn thức $8x^3$ có bậc $3$. \\
b) Thay $x = 1{,}5$ vào $V$, ta được $V = 8 \cdot (1{,}5)^3 = 8 \cdot 3{,}375 = 27\text{ cm}^3$.
	}
\end{bt}

% nguồn: De_AI_Toan8_cuoi_hoc_ki_i_1 (nhập 21/11/2026)
\begin{bt}%[8D1N1-3]%[ma:5b0f99d8]
	Thu gọn đa thức $K = 5x^2y - 3xy + x^2y - 6x^2y + 4xy + 1$ rồi tính giá trị của $K$ tại $x = 3$, $y = -2$.
	\loigiai{
		Ta có $K = (5x^2y + x^2y - 6x^2y) + (-3xy + 4xy) + 1 = xy + 1$. \\
Thay $x = 3$, $y = -2$ vào $K$, ta được $K = 3 \cdot (-2) + 1 = -6 + 1 = -5$.
	}
\end{bt}

% nguồn: De_AI_Toan8_cuoi_hoc_ki_i_1 (nhập 21/11/2026)
\begin{bt}%[8D1N1-3]%[ma:a7d440d2]
	Một cửa hàng bán táo với giá $40$ nghìn đồng/kg và cam với giá $30$ nghìn đồng/kg.
a) Viết đa thức biểu thị số tiền (nghìn đồng) phải trả khi mua $x$ kg táo và $y$ kg cam.
b) Tính số tiền phải trả khi mua $3$ kg táo và $4$ kg cam.
	\loigiai{
		a) Đa thức biểu thị số tiền (nghìn đồng) phải trả là $T = 40x + 30y$. \\
b) Thay $x = 3$, $y = 4$ vào $T$, ta được $T = 40 \cdot 3 + 30 \cdot 4 = 120 + 120 = 240$ (nghìn đồng).
	}
\end{bt}
